\section{Data}
\label{sec:data}

The study uses elevation products, one land cover product, and two radar products that are decoded for the radar availability audit of Section~\ref{sec:results-radar}. Table~\ref{tab:data} lists them. Each
download is recorded with its provider URL, HTTP metadata and SHA-256 digest, and every derived
number in this paper names the artifact it was read from.

\begin{table}[tbp]
\centering
\caption{Data inventory. Grid spacing is the provider's stated posting. The role column says how
each product is used.}
\label{tab:data}
\footnotesize
\begin{tabular}{@{}lp{0.30\linewidth}p{0.12\linewidth}p{0.17\linewidth}p{0.22\linewidth}@{}}
\toprule
Body & Product & Posting & Frame & Role \\
\midrule
Moon & PGDA product 78, Site01 \citep{barker2021} & 5 m & south polar stereographic & target windows, error clones \\
Moon & PGDA product 90, LOLA south polar model, 80 to 90 degrees south & 20 m & south polar stereographic & encoder training \\
Mars & USGS CTX DTMs \citep{laura2021} & 20 m & equirectangular & encoder training, fingerprint distributions \\
Earth & Copernicus GLO-30 \citep{airbus2022copernicus} & 1 arcsec & geographic & search windows, encoder training, probe \\
Earth & ESA WorldCover v200 \citep{zanaga2022} & 10 m & geographic & land cover mask \\
Earth & UAVSAR Holuhraun scene \citep{rosen2006} & multilooked & slant range and ground range & radar audit \\
Moon & Mini-RF south polar CPR mosaic \citep{nozette2010} & mosaic & south polar stereographic & radar audit \\
\bottomrule
\end{tabular}
\end{table}

\paragraph{Moon.} The lunar target is Site01 of the high resolution site products that \citet{barker2021} built from Lunar Orbiter Laser Altimeter ranges \citep{smith2010}, the site that the provider names Connecting ridge. The product gives surface
height at 5 m per pixel in south polar stereographic metres, the number of LOLA returns per pixel,
the total height uncertainty, and \ConfigClones{} error clones of the elevation model. Pixels without a
LOLA return are filled by interpolation in the provider's surface height, so we carry the observed
fraction of each window beside its statistics. The Site01 extent is tiled with \ConfigWindowM\,m
windows at a \ConfigStrideM\,m stride, giving \MoonWindowsGenerated{} windows, of which
\MoonWindowsEligible{} pass the coverage rule. The primary target window is the eligible window
that contains the centroid of the union of the provider's region of interest polygons. Its height
grid is \MoonTargetValidPct\% valid, but for only \MoonTargetObservedPct\% of its valid analysis cells does the nearest 5 m pixel hold a LOLA return; elsewhere the 5 m heights are the provider's interpolation. A permanent shadow fraction is not reported,
because no permanent shadow product was qualified for this window.

\paragraph{Mars.} Martian terrain comes from the USGS Ames Stereo Pipeline CTX elevation models
\citep{laura2021,beyer2018}, built from Context Camera stereo pairs \citep{malin2007}. Each model
ships a quality file with the statistics of its height differences from MOLA altimetry \citep{smith2001} after ground control. Models whose standard deviation of that difference exceeds 10\,m are dropped; for the encoder dataset this left \MarsDtmsUsed{} models and dropped \MarsDtmsDropped. Two Martian target sites are declared, a box at Jezero crater and a box in Arcadia Planitia; each target site is one CTX model that passes the threshold, tiled with the same windows as on the Moon, and Section~\ref{sec:results-mars} gives the rules that fix the target windows.

\paragraph{Earth.} Earth elevation is the Copernicus GLO-30 digital surface model, read as cloud
optimised GeoTIFFs from its public bucket. It is a surface model, so canopy and buildings are part
of the height. We therefore keep only windows whose land cover in ESA WorldCover v200
\citep{zanaga2022} is at least \ConfigMinAllowedPct\% bare or sparse vegetation, snow and ice,
grassland or shrubland, with permanent water excluded. Land cover is sampled as the nearest class
at the centres of the \ConfigGridM\,m analysis cells. Two regions are declared. The Holuhraun lava
field in Iceland is clipped to the footprint of the UAVSAR scene and gives
\HoluhraunWindowsGenerated{} windows, \HoluhraunWindowsEligible{} of them eligible. Rosamond Dry
Lake in California is a terrain only region with \RosamondWindowsGenerated{} windows, of which
\RosamondWindowsEligible{} are eligible and \RosamondWindowsRejected{} fail the land cover rule.
Earth windows are cut in the local UTM zone with the same width and stride as on the Moon.

\paragraph{Windows and resampling.} Every statistic is computed in metres inside each body's own
frame, and lunar, Martian and terrestrial coordinates are never mixed. A window is resampled onto a
square grid about its centre with an inverse spherical azimuthal equidistant projection
\citep{snyder1987}. Pixel centres sit at $x_i = (i + \tfrac12)\Delta - W/2$ and
$y_j = W/2 - (j + \tfrac12)\Delta$ for width $W$ and spacing $\Delta$. The spheres have the mean
radii of the Moon, Mars and Earth (1737400, 3396190 and 6371008.8\,m). Each output
pixel is a bilinear interpolation of its four source neighbours and is nodata if any neighbour is
nodata. No other resampling is used, and the same code runs in Python and in the browser, where a
parity test on shared fixtures holds the two implementations to the tolerance stated in each
fixture. The fingerprint uses a
\ConfigGridM\,m grid; the encoder uses \ConfigInputPx{} by \ConfigInputPx{} pixels at
\ConfigInputSpacing\,m.

\paragraph{Encoder dataset and splits.} Encoder windows are sampled uniformly in projected area from the lunar model within the cap south of 80 degrees, from the Martian models that pass the quality threshold, and from Earth tiles at a fixed number of windows per tile, with the Earth land cover rule applied. They are split by space, never at random: lunar longitude sectors, groups of Martian models, and 2 degree Earth blocks. Windows that straddle a block boundary are dropped, and so is every window within 10\,km of the lunar target or of the Jezero box, in every split; the Arcadia footprint is not set, so no exclusion around it could be checked. The dataset holds \DatasetWindows{} windows. The training split has \DatasetMoonWindows{} lunar, \DatasetMarsWindows{} Martian and \DatasetEarthWindows{} terrestrial windows; validation has \DatasetMoonVal, \DatasetMarsVal{} and \DatasetEarthVal; and test has \DatasetMoonTest, \DatasetMarsTest{} and \DatasetEarthTest; the lunar target window is stored outside the splits. These counts are far below the dataset size planned for the encoder; the dataset was built on a CPU server at a reduced scope, and its small size limits every encoder result.
